% Part: second-order-logic
% Chapter: metatheory
% Section: compactness

\documentclass[../../../include/open-logic-section]{subfiles}

\begin{document}

\olfileid[hi]{sol}{met}{lst}

\olsection{द्वितीय-कोटि तर्कशास्त्र के लिए ल्योवेनहाइम--स्कोलेम प्रमेय विफल होता है}

\begin{explain}
(अवरोही) ल्योवेनहाइम--स्कोलेम प्रमेय कहता है कि !!{sentence} के जिस भी
समुच्चय का कोई अनंत मॉडल है, उसका कोई !!a{enumerable} मॉडल भी है। यह भी
पूर्णता प्रमेय का परिणाम है: पूर्णता का प्रमाण !!{sentence} के किसी भी
संगत समुच्चय के लिए एक मॉडल बनाता है, और वह मॉडल !!{enumerable} होता है।
एक आरोही ल्योवेनहाइम--स्कोलेम प्रमेय भी है, जो आश्वस्त करता है कि यदि
!!{sentence} के किसी समुच्चय का कोई !!a{denumerable} मॉडल है, तो उसका कोई
!!a{nonenumerable} मॉडल भी है। द्वितीय-कोटि तर्कशास्त्र में दोनों प्रमेय
विफल होते हैं।
\end{explain}


\begin{thm}
\ollabel{thm:sol-no-ls} द्वितीय-कोटि तर्कशास्त्र के लिए
ल्योवेनहाइम--स्कोलेम प्रमेय विफल होता है: ऐसे !!{sentence} हैं जिनके अनंत
मॉडल हैं, पर कोई !!{enumerable} मॉडल नहीं है।
\end{thm}

\begin{proof}
स्मरण करें कि
\[
\fn{Count} \ident \lexists[z][\lexists[u][\lforall[X][((X(z) \land
      \lforall[x][(X(x) \lif X(u(x)))]) \lif \lforall[x][X(x)])]]]
\]
!!a{structure}~$\Struct{M}$ में तभी और केवल तभी सत्य है जब $\Domain{M}$
!!{enumerable} हो; अतः $\lnot\fn{Count}$, ~$\Struct{M}$ में तभी और केवल
तभी सत्य है जब $\Domain{M}$ !!{nonenumerable} हो। ऐसी !!{structure}
अस्तित्व में हैं---किसी भी !!{nonenumerable} समुच्चय को !!{domain} मान लें,
जैसे $\Pow{\Nat}$ या~$\Real$। इसलिए $\lnot \fn{Count}$ के अनंत मॉडल तो
हैं, पर कोई !!{enumerable} मॉडल नहीं है।
\end{proof}

\begin{thm}
ऐसे !!{sentence} हैं जिनके !!{denumerable} मॉडल तो हैं, पर कोई
!!{nonenumerable} मॉडल नहीं है।
\end{thm}

\begin{proof}
$\fn{Count} \land \fn{Inf}$, ~$\Nat$ में सत्य है, पर ऐसी किसी
!!{structure}~$\Struct{M}$ में सत्य नहीं है जिसका $\Domain{M}$
!!{nonenumerable} हो।
\end{proof}


\end{document}
